% -----------------------------------------------------------------------
% §8 Implementation
% -----------------------------------------------------------------------
\section{Implementation}
\label{sec:implementation}

This section describes the engineering realisation of PQ-Aliro across two
platforms: an Android Host Card Emulation (HCE) application acting as the User
Device (UD) and a PC-side NFC reader application.
Both implementations are open source and written in Java.

% -----------------------------------------------------------------------
\subsection{System Architecture}
\label{sec:impl:arch}

The UD application runs on Android and exposes the PQ-Aliro applet via the
Host Card Emulation framework (\texttt{HostApduService}).
The Reader application runs on a standard PC equipped with an NFC reader
peripheral and communicates with the UD over the ISO/IEC~14443 Type-A
contactless channel.
Both sides share a common message state machine that advances through four
stages: SELECT $\to$ AUTH0 $\to$ LOADCERT $\to$ AUTH1\@.
Each stage corresponds to a pair of command and response APDUs as detailed in
Section~\ref{sec:apdu-mapping}.

All post-quantum cryptographic operations are provided by BouncyCastle~1.78
with the BCPQC extension.
On the UD, the provider is registered at application start-up:
\begin{verbatim}
  Security.addProvider(new BouncyCastlePQCProvider());
\end{verbatim}
The same provider is used on the PC Reader.
No native (JNI) code is required: both ML-KEM and ML-DSA are implemented
entirely in pure Java within the BouncyCastle library, making the codebase
portable across any JVM-capable platform.
The UD manages long-term Dilithium key material independently of Android
Keystore, since the Android Keystore API does not yet expose lattice-based
algorithm support; all PQC key operations therefore remain in the application
layer, with the Keystore used only to protect symmetric encryption keys (see
Section~\ref{sec:impl:isolation}).

% -----------------------------------------------------------------------
\subsection{ML-KEM Integration}
\label{sec:impl:mlkem}

ML-KEM~\cite{NIST:FIPS203} is realised using the
CRYSTALS-Kyber~\cite{EuroSP:BDKLLSSSS18} implementation in BouncyCastle BCPQC.
The Reader generates a fresh ephemeral key pair for every transaction:

\begin{verbatim}
  KeyPairGenerator kpg =
      KeyPairGenerator.getInstance("Kyber", "BCPQC");
  kpg.initialize(KyberParameterSpec.kyber768, new SecureRandom());
  KeyPair ephemeral = kpg.generateKeyPair();
\end{verbatim}

\noindent
The 1\,184-byte public key is extracted as a raw byte array (the SPKI
outer wrapper is stripped) and transmitted in the AUTH0 Command under TLV
tag \texttt{0x87}.
On the UD, the received bytes are reconstructed into a
\texttt{KyberPublicKeyParameters} object via a helper
\texttt{decodeKyberPublicKey(rawBytes)}, which re-attaches the algorithm
identifier before passing the key to the JCA layer.

Encapsulation on the UD uses a \texttt{KeyGenerator} with a
\texttt{KEMGenerateSpec}:
\begin{verbatim}
  KeyGenerator kg =
      KeyGenerator.getInstance("Kyber", "BCPQC");
  kg.init(new KEMGenerateSpec(readerPublicKey, "AES"));
  SecretKeyWithEncapsulation ske =
      (SecretKeyWithEncapsulation) kg.generateKey();
\end{verbatim}

\noindent
The 1\,088-byte encapsulation $c$ is returned in the AUTH0 Response under
tag \texttt{0x86}.
The 32-byte shared secret $S_e$ is immediately imported into the Android
Keystore under alias \texttt{SSKey}; the raw byte array is then zeroed with
\texttt{Arrays.fill(rawBytes, (byte) 0x00)} to prevent the secret from
persisting on the JVM heap.
The Reader performs decapsulation symmetrically via \texttt{KEMExtractSpec}
to recover $S_e$, after which both parties derive the session key
$\mathit{SK\_device}$ through the KDF chain described in
Section~\ref{sec:design:kdf}.

% -----------------------------------------------------------------------
\subsection{ML-DSA Integration and Key Isolation}
\label{sec:impl:isolation}

ML-DSA~\cite{NIST:FIPS204} is realised from the CRYSTALS-Dilithium
submission~\cite{TCHES:DKLLSSS18} in BouncyCastle BCPQC.

\paragraph{Reader-side signing.}
On the PC Reader, the long-term Dilithium private key is held in process
memory throughout the application lifetime.
Signing uses the standard JCA \texttt{Signature} interface:
\begin{verbatim}
  Signature signer =
      Signature.getInstance("Dilithium", "BCPQC");
  signer.initSign(readerPrivateKey);
  signer.update(transcriptBytes);
  byte[] sigma = signer.sign();
\end{verbatim}

\paragraph{UD-side key isolation.}
The UD cannot safely hold the Dilithium private key in the main application
process.
The root cause is that \texttt{DilithiumPrivateKeyParameters} does not
implement \texttt{javax.security.auth.Destroyable} in BouncyCastle~1.78;
calling \texttt{destroy()} leaves the key material on the JVM heap until the
garbage collector decides to overwrite the memory region, with no timing
guarantee.
A compromised main process could therefore read the private key from heap
memory at any point between key generation and GC collection.

To eliminate this attack surface, the UD delegates all Dilithium signing to a
dedicated \texttt{PQCSignatureService} that is declared with the Android
manifest attribute \texttt{android:isolatedProcess="true"}.
An isolated process runs in a separate Linux process with a distinct address
space, a separate UID, and no file-system access beyond its own
content-provider interface.
The main process sends the message bytes to be signed over a Binder IPC
interface; \texttt{PQCSignatureService} loads the encrypted private key (see
below), performs the signing operation entirely within its own address space,
returns the 3\,309-byte signature over the same Binder channel, and
immediately calls \texttt{Process.killProcess(Process.myPid())} to terminate
itself.
Because the entire lifecycle of the isolated process---from private key
loading to process exit---is bounded to a single signing call, the Dilithium
private key material never appears in the main application's address space and
is forcibly reclaimed by the OS when the isolated process terminates.

\paragraph{Persistent key storage.}
The Dilithium private key cannot be stored in Android Keystore directly.
Instead, it is stored in \texttt{SharedPreferences} encrypted under AES-256-GCM:
the raw private key bytes are encrypted once at provisioning time
(\texttt{KEY\_KYB\_PRIV\_CT} stores the ciphertext, \texttt{KEY\_KYB\_PRIV\_IV}
stores the GCM IV), and the 256-bit AES wrapping key is held in Android
Keystore, which provides hardware-backed protection on devices with a Trusted
Execution Environment.
The isolated process retrieves and decrypts the private key on each invocation,
uses it exactly once, then discards it along with the process.

% -----------------------------------------------------------------------
\subsection{NFC Fragmentation and Asynchronous Signing}
\label{sec:impl:frag}

\paragraph{APDU fragmentation.}
As described in Section~\ref{sec:apdu-mapping}, PQ-Aliro payloads must be
split into chains of APDUs to respect the ISO/IEC~14443 ISO-DEP frame limit.
The Reader sends AUTH0 and AUTH1 command chunks of at most 200 bytes, and
LOADCERT chunks of at most 255 bytes.
The UD sends response chunks of at most 240 bytes.
Fields longer than 255 bytes are encoded with the BER-TLV long-form length
prefix \texttt{0x82\,\textit{len\_hi}\,\textit{len\_lo}}, adding 4 bytes of
overhead per affected field.
This encoding is applied to the ML-KEM-768 ephemeral public key (tag
\texttt{0x87}), the KEM ciphertext (tag \texttt{0x86}), all ML-DSA-65
public keys (tags \texttt{0x5A}), and all ML-DSA-65 signatures (tag
\texttt{0x9E}) in both AUTH1 Command and AUTH1 Response.

\paragraph{Asynchronous signing to prevent HCE deadlock.}
Android's HCE framework invokes \texttt{processCommandApdu()} on a dedicated
system thread.
If this callback blocks on a synchronous Binder call to an isolated process,
the system-thread call stack is held indefinitely: the Reader's NFC controller
times out the APDU exchange, and Android may trigger an Application Not
Responding (ANR) dialog.
To avoid this, the UD adopts a non-blocking response strategy for the AUTH1
Command.
Upon receiving the final chunk of AUTH1, \texttt{processCommandApdu()} spawns
a background thread to perform the Binder-based signing call and immediately
returns the two-byte response \texttt{0x61\,0x00} (``more data available, zero
bytes ready'') to the NFC system thread, releasing it at once.
The Reader interprets \texttt{0x61\,\textit{xx}} as a continuation indicator
and enters a polling loop, sending \texttt{GET RESPONSE} (\texttt{INS = 0xC0})
commands at regular intervals.
Once the background signing thread completes and writes the fully assembled
AUTH1 Response into the \texttt{responseQueue}, the next \texttt{GET RESPONSE}
retrieves the first 240-byte chunk, with subsequent \texttt{GET RESPONSE}
commands collecting the remaining 22 chunks.

\paragraph{AEAD and IV management.}
The AUTH1 Response payload is protected by AES-256-GCM using the session key
$\mathit{SK\_device}$.
The GCM IV is derived as the 12-byte value
$\texttt{0x00\,00\,00\,00\,00\,00\,00\,01} \mathbin{\|} \mathit{ctr}$,
where $\mathit{ctr}$ is a 4-byte big-endian counter initialised to 1 and
incremented for each subsequent AEAD invocation within a session.
The IV is not transmitted over NFC; both parties derive it independently from
the shared counter state.
The 16-byte GCM authentication tag is appended to the ciphertext, yielding
a 5\,289-byte AUTH1 Response payload.
All citation to the Aliro~1.0 specification~\cite{CSA:Aliro10} for the
APDU command structure and status-word conventions has been preserved
throughout this implementation.

% -----------------------------------------------------------------------
\subsection{End-to-End Testing}
\label{sec:impl:testing}

The implementation is validated by a comprehensive end-to-end test suite
that exercises all nine PQC parameter combinations:
Dilithium$\{2,3,5\} \times$ Kyber$\{512,768,1024\}$.
The test harness runs entirely in-memory: the PC-side Reader and the Android
HCE logic are instantiated within the same JVM process, with a byte-array
channel substituting for the physical NFC interface.
This removes the need for real NFC hardware and makes the tests fully
reproducible in a standard CI environment.

Each test configuration exercises 20 assertions covering: (i) KEM
correctness---that the shared secret $S_e$ recovered by the Reader via
\texttt{KEMExtractSpec} matches the value encapsulated by the UD; (ii)
transcript hash consistency---that $\mathit{transAH0} =
\mathrm{SHAKE256\text{-}512}(\mathit{AUTH0\_cmd} \mathbin{\|} \mathit{AUTH0\_resp})$
is identical on both sides; (iii) KDF consistency---that
$\mathit{SK\_device}$ derived by both parties is identical;
(iv) AEAD round-trip correctness; (v) ML-DSA signature verification in both
directions (Reader signature verified by UD, UD signature verified by Reader);
and (vi) serialisation integrity of all TLV fields across the fragmentation
and reassembly pipeline.
All 20 assertions pass across all nine parameter combinations, confirming
the correctness of the implementation for the full range of NIST security
levels supported by PQ-Aliro~\cite{NIST3PQC:Lima21}.


